\begin{table}[tb]
\centering\small\setlength{\tabcolsep}{4pt}
\caption{全４課題のH−M費用差（論理token）．反復２以外は順序対応と固定推論器の計数に基づく区間であり，統計的信頼区間ではない．元の不明生成は残す．}
\begin{tabular}{lr}
集計 & H−M \\ \hline
反復1 & +32,959〜+34,925 \\
反復2 & +69,581 \\
反復3 & -127,288〜-123,560 \\
合計 & -24,748〜-19,054 \\
\hline
\end{tabular}
\end{table}
事後の補助解析では，P80とP81の所有推論器の完全な直列ログを全429生成へ順序対応させた．既知427件すべてで，入力処理数＝終端入力−cache入力，生成数＝終端出力，送信前計数＝終端入力が一致した．
不明要求$j$について，同一送信の入力計数を$\widehat I_j$，ログに残った最大生成進捗を$g_j$，要求の出力上限を$B_j$とすると，固定推論器の計数・停止実装を前提に
\begin{equation}\widehat I_j+g_j\leq T_j\leq\widehat I_j+B_j\end{equation}
を得る．生成数は入力処理完了後に増加し，要求上限で停止する．全保存要求のhashと上限を照合し，別名による上書きがないことを確認した．不明２件の区間は順に39,657〜41,623，40,971〜44,699 tokenである．
この順序対応は既知427件の一致で裏付けられるが，要求hashとサーバtask IDを暗号学的に結合したものではない．記録外の生成がないこと，入力計数と生成が同じmessages/toolsを使うこと，固定実装どおりに処理されることを前提とする．終端使用量の補完ではなく，cache入力と取消し後の処理を含む論理tokenの範囲である．GPU演算量，投機生成の未受理作業，料金や実時間の区間ではない．
３反復合計ではHが条件付きで19,054〜24,748 token，約0.79〜1.02\%少ない．ただし両条件で成功した３課題の減少397,675 tokenに対し，両条件で失敗した書出し課題は372,927〜378,621 token増えている．反復ごとの全費用差も増加・増加・減少であり，安定した改善とはいえない．第３反復の書出し費用差は双方の未達で停止した差であり，完成に必要な費用の削減ではない．元の不明２件と採用基準未達は保持する．
